% =============================================================================
%  PREÁMBULO DE ESTILO  —  versión corregida
%  Cambios respecto al original:
%    1. glosstex  →  glossaries-extra  (sistema único de acrónimos)
%    2. subfig    →  subcaption        (paquete moderno, no deprecado)
%    3. \ifpdf/hyperref  →  carga directa moderna
%    4. Geometría TikZ hardcodeada  →  relativa a \textwidth/\linewidth
%    5. Código comentado de capítulos eliminado
% =============================================================================

% -----------------------------------------------------------------------------
% TIPOGRAFÍA E IDIOMA
% -----------------------------------------------------------------------------
\usepackage[T1]{fontenc}
\usepackage[utf8]{inputenc}
\usepackage[english, spanish, es-noindentfirst, activeacute]{babel}

% Sobreescribimos los nombres que babel asigna en español
% ("cuadro" → "Tabla", etc.) para convenciones latinoamericanas.
\addto\captionsspanish{%
  \def\tablename{Tabla}%
  \def\listtablename{\'{I}ndice de Tablas}%
  \def\contentsname{\'{I}ndice}%
  \def\chaptername{Cap\'{i}tulo}%
}

% -----------------------------------------------------------------------------
% HIPERVÍNCULOS  (carga directa — \ifpdf es innecesario con motores modernos:
%                 pdflatex, xelatex y lualatex producen PDF de forma nativa)
% -----------------------------------------------------------------------------
\usepackage[
  pdfborder  = {0 0 0},
  colorlinks = true,
  linkcolor  = purple,
  filecolor  = purple,
  citecolor  = purple,
  urlcolor   = purple,
  hyperindex,
  pdfpagelabels
]{hyperref}
\usepackage{bookmark}   % debe ir después de hyperref

% -----------------------------------------------------------------------------
% ACRÓNIMOS Y GLOSARIOS  —  glossaries-extra (sistema único)
%
%   • Se elimina el bloque \ifx\generaacronimos que cargaba el paquete
%     obsoleto "glosstex" y dejaba comandos vacíos como fallback.
%   • glossaries-extra cubre todos los casos: si no se usan entradas
%     simplemente no se imprime nada.
%   • La opción [abbreviations] es el mecanismo nativo de glossaries-extra
%     para siglas; [acronym] también funciona si usas \newacronym.
%   • \makeglossaries DEBE ir después de \usepackage{glossaries-extra}.
% -----------------------------------------------------------------------------
\usepackage[
  toc,           % añade los glosarios al índice general
  acronym,       % crea el glosario tipo \acronymtype (compatible con \newacronym)
  % abbreviations % alternativa nativa de glossaries-extra (usa \newabbreviation)
]{glossaries-extra}
\makeglossaries   % OBLIGATORIO aquí, genera los archivos .acn/.glo

% -----------------------------------------------------------------------------
% ÍNDICE GENERAL
% -----------------------------------------------------------------------------
\usepackage{makeidx}

\ifx\generaindice\undefined
\else
  \makeindex
\fi

% -----------------------------------------------------------------------------
% FIGURAS Y SUBFIGURAS
%   subfig está deprecado; subcaption es el reemplazo moderno y es compatible
%   con memoir, KOMA-Script y caption.
% -----------------------------------------------------------------------------
\usepackage{graphicx}
\usepackage{epsfig}
\usepackage[
  font       = small,
  labelfont  = bf,
  format     = plain,
  justification = centering
]{caption}
\usepackage{subcaption}   % reemplaza a subfig

% Entorno SubFloat reimplementado sobre subcaption.
% Uso:
%   \begin{SubFloat}[etiqueta TOC]{Leyenda visible}
%     <contenido>
%   \end{SubFloat}
%
% Si #1 está vacío, usa solo la leyenda visible como etiqueta del subfloat.
\makeatletter
\newbox\sf@box
\newenvironment{SubFloat}[2][]{%
  \def\sf@one{#1}%
  \def\sf@two{#2}%
  \setbox\sf@box\hbox\bgroup
}{%
  \egroup
  \ifx\@empty\sf@one\@empty\relax
    \begin{subfigure}[t]{\wd\sf@box}%
      \usebox\sf@box
      \caption{\sf@two}%
    \end{subfigure}%
  \else
    \begin{subfigure}[t]{\wd\sf@box}%
      \usebox\sf@box
      \caption[\sf@one]{\sf@two}%
    \end{subfigure}%
  \fi
}
\makeatother

% -----------------------------------------------------------------------------
% TABLAS
% -----------------------------------------------------------------------------
\usepackage{tabularx}
\usepackage{multicol}

% -----------------------------------------------------------------------------
% LISTADOS DE CÓDIGO
% -----------------------------------------------------------------------------
\usepackage{listings}
\lstset{
  basicstyle      = \small\ttfamily,
  showstringspaces= false,
  tabsize         = 3
}

% -----------------------------------------------------------------------------
% PAQUETES DE UTILIDAD
% -----------------------------------------------------------------------------
\usepackage{verbatim}
\usepackage{shortvrb}
\usepackage{url}
\usepackage{calc}
\usepackage{xcolor}
\usepackage{tikz}
\usetikzlibrary{shapes.misc}

% Espaciado entre párrafos
\setlength{\parskip}{0.2ex}

% -----------------------------------------------------------------------------
% FORMATO DE CAPÍTULOS
% -----------------------------------------------------------------------------
\usepackage{titlesec}

\titleformat{\chapter}[display]
  {\vfill\centering\normalfont}
  {%
    \fontsize{48pt}{48pt}\usefont{T1}{phv}{m}{n}%
    \MakeUppercase{\chaptername}%
    \quad
    {\fontsize{48pt}{48pt}\selectfont\thechapter}%
  }
  {0pt}
  {%
    \vspace{4pt}%
    \hrule height 0.8pt%
    \vspace{8pt}%
    \Huge\usefont{T1}{phv}{b}{n}\MakeUppercase
  }
  [\vspace{6pt}\hrule height 0.4pt\vfill\clearpage]
\titlespacing*{\chapter}{0pt}{0pt}{0pt}

% -----------------------------------------------------------------------------
% FORMATO DE SECCIONES  —  APA 7, hoja carta / A4
%
%  REQUISITO: barra y nodo comienzan en el sangrado APA 7 de 1 pulgada
%  (2.54 cm); el título de sección queda justo después del nodo.
%
%  ESTRATEGIA (corregida):
%   1. \titlespacing* provee el sangrado izquierdo de 1 pulgada.
%      x=0 dentro del tikzpicture corresponde al borde del sangrado.
%
%   2. La barra se dibuja con \begin{scope}[overlay]...\end{scope}.
%      La clave "overlay" excluye ese contenido del bounding-box del
%      tikzpicture, por lo que NO altera el ancho del label en \titleformat.
%
%   3. El nodo (rounded rectangle) se dibuja fuera del scope overlay:
%      su ancho SÍ contribuye al bounding-box. Por tanto el label box
%      tiene exactamente el ancho del nodo.
%
%   4. \titleformat coloca el título a {sep} del borde derecho del label,
%      es decir, inmediatamente después del nodo. ✓
%
%   5. Sin trim left / trim right: ya no son necesarios.
%
%  DIAGNÓSTICO del bug anterior:
%      trim right = \sectionbarwidth - \nodelabelwidth se calculaba en el
%      preámbulo cuando \textwidth aún podía valer 0pt (antes de que
%      geometry lo fijara), resultando en trim right ≈ −\nodelabelwidth.
%      Con ese valor negativo el label box era MÁS ANCHO que la barra,
%      el título quedaba al extremo derecho o se salía del margen.
% -----------------------------------------------------------------------------

% ── Parámetros APA 7 ─────────────────────────────────────────────────────────
\newlength{\APAindent}
\setlength{\APAindent}{2.54cm}   % 1 pulgada exacta

% ── \section ─────────────────────────────────────────────────────────────────
\newcommand\titlebar{%
  \tikz[baseline=(nd.base)]{%
    % Barra de fondo — overlay: no altera el bounding-box del label
    \begin{scope}[overlay]
      \fill[cyan!25]
        (0pt, -1ex) rectangle (\dimexpr\linewidth\relax, 2.5ex);
    \end{scope}
    % Nodo — SÍ contribuye al bounding-box → define el ancho del label
    \node[
      fill           = cyan!60!white,
      anchor         = base west,
      rounded rectangle,
      minimum height = 3.5ex,
      inner sep      = 5pt
    ] (nd) at (0pt, 0) {\textbf{\thesection.}};
  }%
}

\titleformat{\section}
  {\large}{\titlebar}{0.6em}{}

% left = \APAindent aplica el sangrado APA 7.
% Dentro del tikzpicture, \linewidth = \textwidth − \APAindent (ajustado por LaTeX).
\titlespacing*{\section}{\APAindent}{1.5ex plus .4ex minus .2ex}{0.8ex}

% ── \subsection ──────────────────────────────────────────────────────────────
\newcommand{\subtitlebar}{%
  \tikz[baseline=(nd.base)]{%
    % Barra delgada — overlay
    \begin{scope}[overlay]
      \fill[cyan!25]
        (0pt, -0.5ex) rectangle (\dimexpr\linewidth\relax, 1.8ex);
    \end{scope}
    % Nodo
    \node[
      fill           = cyan!60!white,
      anchor         = base west,
      rounded rectangle,
      minimum height = 2.2ex,
      inner sep      = 4pt
    ] (nd) at (0pt, 0) {\textbf{\thesubsection.}};
  }%
}

\titleformat{\subsection}
  {\large}{\subtitlebar}{0.6em}{}

\titlespacing*{\subsection}{\APAindent}{1.2ex plus .3ex minus .2ex}{0.6ex}
